\section*{Capítulo \ref{ch:b2:huckel} --- Teoría de Hückel y sistemas conjugados}

\begin{solution}{exo:b2:huckel:1}
Niveles del alilo ($n = 3$): $\alpha + 1.414\beta$, $\alpha$, $\alpha - 1.414\beta$. Catión (2 electrones):
$2\alpha + 2.828\beta$; radical (3): $3\alpha + 2.828\beta$; anión (4): $4\alpha + 2.828\beta$. Los electrones
adicionales van al nivel no enlazante y solo cambian la parte en $\alpha$.
\end{solution}

\begin{solution}{exo:b2:huckel:2}
$2(\alpha + 1.618\beta) + 2(\alpha + 0.618\beta) = 4\alpha + 4.472\beta$.
\end{solution}

\begin{solution}{exo:b2:huckel:3}
Benceno (6), anión ciclopentadienilo (6) y tropilio (6): \omterm{def:b2:huckel:aromatic}{aromáticos}. Catión
ciclopentadienilo (4), ciclobutadieno (4) y ciclooctatetraeno plano (8): \omterm{def:b2:huckel:aromatic}{antiaromáticos} si son planos
(el ciclooctatetraeno real tiene forma de bañera y es no \omterm{def:b2:huckel:aromatic}{aromático}).
\end{solution}

\begin{solution}{exo:b2:huckel:4}
Un octógono en el círculo: niveles $\alpha + 2\beta$, $\alpha + 1.414\beta$ (dos veces), $\alpha$ (dos veces),
$\alpha - 1.414\beta$ (dos veces), $\alpha - 2\beta$. Ocho electrones: dos en el más bajo, cuatro en la
primera pareja y uno en cada orbital de la pareja no enlazante.
\end{solution}

\begin{solution}{exo:b2:huckel:5}
$4.472\beta - 4\beta = 0.472\beta$, es decir, $0.472 \times 75 = \qty{35}{kJ/mol}$, frente a unos
\qty{10}{kJ/mol} según los datos de hidrogenación: el modelo de Hückel, con $\beta$ ajustado al benceno,
sobrestima la conjugación de las cadenas abiertas.
\end{solution}

\begin{solution}{exo:b2:huckel:6}
$p_{12} = p_{34} = 0.894$, $p_{23} = 0.447$: los enlaces de los extremos, C1--C2 y C3--C4, son los
más cortos, y el enlace central, más largo pero más corto que un enlace sencillo.
\end{solution}

\begin{solution}{exo:b2:huckel:7}
Anión, seis electrones: $2 \times 2 + 4 \times 0.618$, $E_\pi = 6\alpha + 6.472\beta$, una capa cerrada:
\omterm{def:b2:huckel:aromatic}{aromático}. Catión, cuatro: $2 \times 2 + 2 \times 0.618$, $4\alpha + 5.236\beta$, con un electrón
en cada orbital de la pareja degenerada: \omterm{def:b2:huckel:aromatic}{antiaromático}, muy inestable.
\end{solution}

\begin{solution}{exo:b2:huckel:8}
Anillo de siete eslabones: $\alpha + 2\beta$, $\alpha + 1.247\beta$ (dos veces) y luego niveles antienlazantes. Seis
electrones, como en el catión \ce{C7H7+}, llenan exactamente los niveles enlazantes: un catión
\omterm{def:b2:huckel:aromatic}{aromático}. El compuesto se ioniza con facilidad en tropilio y bromuro.
\end{solution}

\begin{solution}{exo:b2:huckel:9}
$1.802 + 1.247 + 0.445 - 0.445 - 1.247 - 1.802 = 0$: la suma de las energías es $6\alpha$.
\end{solution}

\begin{solution}{exo:b2:huckel:10}
$n = 6$: $m_k = 2\cos(k\pi/7) = 1.802$, $1.247$, $0.445$, $-0.445$, $-1.247$, $-1.802$. Seis
electrones: $E_\pi = 6\alpha + 2(1.802 + 1.247 + 0.445)\beta = 6\alpha + 6.988\beta$;
deslocalización, $0.988\beta$.
\end{solution}

\begin{solution}{exo:b2:huckel:11}
En el cuadrado, los dos electrones de mayor energía están uno en cada orbital de una
pareja degenerada no enlazante. Alargar dos enlaces opuestos levanta la degeneración:
un orbital baja y el otro sube, y ambos electrones se aparean en el más bajo, lo que
rebaja la energía. La molécula se estabiliza como un rectángulo con dos dobles enlaces
localizados; sigue siendo \omterm{def:b2:huckel:aromatic}{antiaromática} y extremadamente reactiva.
\end{solution}

\begin{solution}{exo:b2:huckel:12}
Matriz (en unidades de $\beta$, con $\alpha$ como origen) $\begin{pmatrix}0 & 1\\ 1 & 1\end{pmatrix}$:
$m = 1.618$ y $-0.618$, $E = \alpha + 1.618\beta$ (enlazante) y $\alpha - 0.618\beta$. Para el
\omterm{def:b2:diatomic-mos:bonding}{orbital enlazante} $c_X = 1.618c_C$, luego $c_C^2 = 0.276$, $c_X^2 = 0.724$; con dos electrones,
$q_C = 0.55$, $q_X = 1.45$: la densidad $\pi$ es atraída hacia X.
\end{solution}

\begin{solution}{pb:b2:huckel:1}
\textbf{1.} $-123.1 - (-4.3) = \qty{-118.8}{kJ/mol}$.
\textbf{2.} $-123.1 - 104.6 = \qty{-227.7}{kJ/mol}$.
\textbf{3.} $-123.1 - 82.9 = \qty{-206.0}{kJ/mol}$.
\textbf{4.} $3 \times (-118.8) = \qty{-356.4}{kJ/mol}$.
\textbf{5.} $356.4 - 206.0 = \qty{150}{kJ/mol}$.
\textbf{6.} El dieno libera $2 \times 118.8 - 227.7 = \qty{10}{kJ/mol}$ menos que dos
dobles enlaces aislados: la conjugación de una cadena abierta da una ganancia pequeña, y el anillo
cerrado del benceno, quince veces más.
\textbf{7.} Una matriz $6 \times 6$ con 1 entre cada átomo y sus dos vecinos (1--2, 2--3,
\dots, 6--1) y 0 en el resto.
\textbf{8.} $m = 2\cos(2\pi k/6)$: $2$, $1$, $1$, $-1$, $-1$, $-2$: $E = \alpha + 2\beta$, $\alpha + \beta$ (dos veces),
$\alpha - \beta$ (dos veces), $\alpha - 2\beta$.
\textbf{9.} Dos electrones en $\alpha + 2\beta$ y cuatro en la pareja $\alpha + \beta$.
\textbf{10.} $E_\pi = 6\alpha + 8\beta$.
\textbf{11.} $3(2\alpha + 2\beta) = 6\alpha + 6\beta$.
\textbf{12.} $2\beta$, es decir, $2|\beta|$ de estabilización.
\textbf{13.} $2 + 1 + 1 - 1 - 1 - 2 = 0$: los seis niveles suman $6\alpha$.
\textbf{14.} $2|\beta| = \qty{150}{kJ/mol}$: $|\beta| = \qty{75}{kJ/mol}$.
\textbf{15.} $0.472 \times 75 = \qty{35}{kJ/mol}$, frente a \qty{10}{kJ/mol}: el modelo ajustado
al benceno sobrestima la cadena abierta.
\textbf{16.} Los orbitales ocupados pueden tomarse como $\psi_1 = (1, 1, 1, 1, 1, 1)/\sqrt6$,
$\psi_2 = (2, 1, -1, -2, -1, 1)/\sqrt{12}$ y $\psi_3 = (0, 1, 1, 0, -1, -1)/2$. Para el enlace
1--2: $p_{12} = 2(\frac16 + \frac{2}{12} + 0) = \frac23$; por simetría, todos los enlaces tienen $p = 2/3$.
\textbf{17.} Entre los de los enlaces de los extremos (0.894) y del enlace central (0.447) del butadieno.
\textbf{18.} Los seis enlaces son equivalentes por simetría; su longitud, \qty{139.7}{pm}, está
entre el doble enlace del eteno (\qty{133.9}{pm}) y el enlace sencillo del etano
(\qty{153.6}{pm}).
\textbf{19.} $0.988|\beta| = \qty{74}{kJ/mol}$, la mitad que la del benceno: cerrar el anillo duplica la
\omterm{def:b2:huckel:delocalisation}{energía de deslocalización}.
\textbf{20.} $\alpha + 2\beta$, $\alpha + 1.414\beta$ (dos veces), $\alpha$ (dos veces), $\alpha - 1.414\beta$ (dos veces),
$\alpha - 2\beta$.
\textbf{21.} Los dos últimos electrones van uno a cada orbital de la pareja no enlazante: una capa
abierta, \omterm{def:b2:huckel:aromatic}{antiaromática}.
\textbf{22.} Se pliega en forma de bañera: los orbitales p de los dobles enlaces vecinos ya no son
paralelos, los enlaces alternan y la molécula se comporta como un polieno.
\textbf{23.} Cuatro electrones: la misma pareja llena a medias, peor aún porque la molécula no puede
evitar la planaridad; se deforma en un rectángulo y es extremadamente reactiva.
\textbf{24.} Un sistema conjugado monocíclico plano es \omterm{def:b2:huckel:aromatic}{aromático} con $4N + 2$ electrones $\pi$,
y \omterm{def:b2:huckel:aromatic}{antiaromático} con $4N$.
\textbf{25.} $\boldsymbol{|\beta| \approx \qty{75}{kJ/mol}}$.
\end{solution}
